% Hodina 30, 1. ročník SPŠ. Rozsah: Rozpis_vyuky_MAT_1_rocnik.tex, řádek 30.

% Plán: 1_rocnik/tematicke_plany/SPS.pdf; společná projektová šablona.

% Učebnice: ucebnice/01_Zakladni_poznatky/Mocniny_a_odmocniny/Mocniny a Odmocniny.tex

% Původ úloh: P1: učebnice, „Odmocniny“. Ostatní vlastní zadání; metodické pořadí a odhad inspirovány Realisticky.cz.

% Metodický zdroj: externi_zdroje/realisticky/Odmocniny/07_Druha_odmocnina.pdf

% URL: http://www.realisticky.cz/ucebnice/01%20Matematika%20S%C5%A0/01%20Z%C3%A1kladn%C3%AD%20poznatky/02%20%C4%8C%C3%ADseln%C3%A9%20obory/07%20Druh%C3%A1%20odmocnina.pdf

% Metodický zdroj: externi_zdroje/realisticky/Odmocniny/08_Treti_odmocnina.pdf

% URL: http://www.realisticky.cz/ucebnice/01%20Matematika%20S%C5%A0/01%20Z%C3%A1kladn%C3%AD%20poznatky/02%20%C4%8C%C3%ADseln%C3%A9%20obory/08%20T%C5%99et%C3%AD%20odmocnina.pdf

\documentclass[11pt,a4paper]{article}

% Spolecny styl pro ucitelske pripravy MAT (Work/MAT)
% Verze 1.1, 2026-10-23
\NeedsTeXFormat{LaTeX2e}
\ProvidesPackage{priprava}[2026/10/23 v1.1 Sablona ucitelskych priprav MAT]

% Zakladni typografie a cestina
\RequirePackage[T1]{fontenc}
\RequirePackage[utf8]{inputenc}
\RequirePackage[czech]{babel}
\RequirePackage{lmodern}
\RequirePackage{microtype}
\RequirePackage{amsmath,amssymb}
\RequirePackage{xcolor}
\RequirePackage{enumitem}
\RequirePackage[most]{tcolorbox}
\RequirePackage{tabularx,booktabs,array}
\RequirePackage{fancyhdr}
\RequirePackage{needspace}
\RequirePackage{geometry}
\RequirePackage[hidelinks]{hyperref}

\geometry{a4paper,top=17mm,bottom=18mm,left=23mm,right=23mm}
\definecolor{notegray}{gray}{0.39}
\setlength{\parindent}{0pt}
\setlength{\parskip}{2.5pt plus 1pt minus 0.5pt}
\setlist[itemize]{topsep=2pt,itemsep=1pt,parsep=0pt,leftmargin=1.45em}
\setlist[enumerate]{topsep=2pt,itemsep=2pt,parsep=0pt,leftmargin=1.65em}
\setcounter{secnumdepth}{0}
\renewcommand{\arraystretch}{1.12}

% Zahlavi: vlevo rocnik | globalni tema, vpravo cislo: nazev hodiny
% Zapati zustava zcela prazdne.
\pagestyle{fancy}
\fancyhf{}
\renewcommand{\headrulewidth}{0.35pt}
\renewcommand{\footrulewidth}{0pt}
\setlength{\headheight}{14pt}
\newcommand{\nastavPripravu}[4]{%
  \fancyhead[L]{\small #1\enspace\textbar\enspace #2}%
  \fancyhead[R]{\small #3: #4}%
}

% Minimalisticke cernobile boxy. Kazda definice zvlast.
\tcbset{commonlesson/.style={enhanced,breakable=false,colback=white,colframe=black,
  colbacktitle=white,coltitle=black,boxrule=0.45pt,
  arc=0.8mm,outer arc=0.8mm,left=3mm,right=3mm,top=1.6mm,bottom=1.6mm,
  toptitle=1mm,bottomtitle=0.7mm,fonttitle=\bfseries,
  attach boxed title to top left={xshift=2.5mm,yshift=-1.15mm},
  boxed title style={colback=white,colframe=white,boxrule=0pt,left=1mm,right=1mm}}}
\newtcolorbox{zadaniBox}[1][]{commonlesson,title={Zadání},#1}
\newtcolorbox{prikladBox}[1][]{commonlesson,title={Příklad},#1}
\newtcolorbox{definiceBox}[1][]{commonlesson,title={Definice},#1}

% Spolecny box pro zadani a jeho reseni, bez mezery mezi castmi.
% Pouziti:
%   \begin{prikladReseniBox}
%     Zadani ulohy.
%     \tcblower
%     \begin{align*} ... \end{align*}
%   \end{prikladReseniBox}
% Spodni cast ma jemny sedy podklad a delici linku; lze v ni sazet
% vice radku matematiky. Obe casti maji jen jeden spolecny ramecek.
\tcbset{solutionlesson/.style={
  commonlesson,
  enhanced jigsaw,
  breakable,
  colbacklower=black!3,
  segmentation style={solid,draw=black!25,line width=0.35pt},
  fontlower=\small\color{notegray},
  before lower={\noindent\textbf{Řešení:}\enspace},
  leftlower=3mm,rightlower=3mm
}}
\newtcolorbox{prikladReseniBox}[1][]{solutionlesson,title={Příklad},#1}
\newtcolorbox{zadaniReseniBox}[1][]{solutionlesson,title={Zadání},#1}

% Pomucky pro autora pripravy: sedive poznamky; \reseni pro starsi pripravy
\newcommand{\poznamka}[1]{\par\begingroup\small\color{notegray}\itshape
  \textnormal{\textbf{Poznámka pro učitele.}}\enspace #1\par\endgroup}
\newcommand{\reseni}[1]{\par\begingroup\small\color{notegray}
  \textbf{Očekávané řešení:}\enspace #1\par\endgroup}
\newcommand{\blok}[1]{\needspace{5\baselineskip}\subsection*{#1}}
\newcommand{\vyrok}[1]{\uv{#1}}


\nastavPripravu{1. ročník}{Odmocniny}{30}{Druhá a třetí odmocnina}

\begin{document}

\blok{Výklad a souvislosti}

Ze známého obsahu čtverce hledáme délku jeho strany. Požadavek nezáporné délky vysvětluje, proč zápis $\sqrt9$ označuje pouze $3$. Třetí mocnina naproti tomu rozlišuje kladná a záporná čísla i bez omezení znaménka.

\begin{definiceBox}[title={Definice},breakable=false]
\textbf{Druhá odmocnina} z $a\ge0$ je jediné nezáporné číslo $b$, pro které platí $b^2=a$. Zapisujeme $b=\sqrt a$. Pro $a<0$ není druhá odmocnina v $\mathbb R$ definována.
\end{definiceBox}


\begin{definiceBox}[title={Definice},breakable=false]
\textbf{Třetí odmocnina} z reálného čísla $a$ je jediné reálné číslo $b$, pro které platí $b^3=a$. Zapisujeme $b=\sqrt[3]a$.
\end{definiceBox}


Každý výsledek ověř příslušným umocněním a u druhé odmocniny také znaménkem. Zatím nepoužíváme pravidla pro odmocninu součinu. Odhad získáme porovnáním se známými druhými nebo třetími mocninami.

\blok{Řešené příklady}

\begin{prikladReseniBox}[title={Příklad 1},breakable=false]
Urči $\sqrt{25}$, $\sqrt[3]{27}$, $\sqrt[3]{-8}$ a $\sqrt{-9}$.
\tcblower
Postupně $5$, $3$, $-2$; poslední výraz nemá v $\mathbb R$ smysl. Kontroly: $5^2=25$, $3^3=27$, $(-2)^3=-8$.
\end{prikladReseniBox}


\begin{prikladReseniBox}[title={Příklad 2},breakable=false]
Porovnej $\sqrt{16}$ s čísly, jejichž druhá mocnina je $16$. Odhadni $\sqrt{39}$.
\tcblower
$\sqrt{16}=4$. Druhou mocninu $16$ mají $4$ i $-4$, ale odmocnina je nezáporná. Protože $6^2<39<7^2$, platí $6<\sqrt{39}<7$.
\end{prikladReseniBox}


\blok{Samostatné procvičení}

\begin{zadaniReseniBox}[title={Procvičení},breakable=false]
\begin{enumerate}[label=\alph*)]
\item Vypočítej $\sqrt{0}$, $\sqrt{0{,}09}$ a $\sqrt[3]{0{,}001}$.
\item Vypočítej $\sqrt[3]{-125}$ a $\sqrt{(-7)^2}$.
\item Najdi dvě sousední celá čísla, mezi nimiž leží $\sqrt7$.
\end{enumerate}
\tcblower
\begin{enumerate}[label=\alph*)]
\item $0$, $0{,}3$ a $0{,}1$; ověříme umocněním.
\item $-5$ a $7$; druhá odmocnina vrací nezáporné číslo.
\item $2<\sqrt7<3$, protože $4<7<9$.
\end{enumerate}
\end{zadaniReseniBox}


\poznamka{Nepiš $\sqrt{16}=\pm4$. Nezaměňuj existenci přesné hodnoty odmocniny s možností vyjádřit ji konečným desetinným zápisem. Iracionální číslo má neperiodický desetinný rozvoj. Tuto skutečnost připomeň, ale nevyžaduj nový důkaz.}

\end{document}
